% Section 06 - Placement QR sur Documents
\section{Placement avancé de QR code sur documents}

\subsection{Vue d'ensemble}

Le menu \menu{Placement QR sur Documents} offre une fonctionnalité similaire à \menu{QR Codes sur PDF}, mais avec deux modes de fonctionnement : manuel et Excel.

\info{Cette fonctionnalité est DISPONIBLE même en mode démo, contrairement à "QR Codes sur PDF".}

\textbf{Principales différences avec "QR Codes sur PDF" :}

\begin{table}[h]
\centering
\begin{tabular}{|l|p{5cm}|p{5cm}|}
\hline
\textbf{Critère} & \textbf{QR Codes sur PDF} & \textbf{Placement QR sur Documents} \\
\hline
Mode démo & ✗ Désactivé & ✓ Disponible \\
\hline
Modes & Uniquement Excel & Manuel ET Excel \\
\hline
QR personnalisé & Basé sur colonnes Excel & Texte libre + Excel \\
\hline
Complexité & Simple & Avancée \\
\hline
\end{tabular}
\caption{Comparaison des deux menus QR}
\end{table}

\subsection{Mode 1 : Placement manuel (texte libre)}

\subsubsection{Cas d'usage}

Utilisez ce mode pour :
\begin{itemize}[leftmargin=*]
    \item Ajouter un QR code unique sur un seul document
    \item Créer un QR code avec du texte personnalisé (URL, vCard, texte libre)
    \item Tester rapidement le positionnement avant un traitement en masse
\end{itemize}

\subsubsection{Workflow manuel}

\begin{enumerate}[leftmargin=*]
    \item Cliquez sur \menu{Placement QR sur Documents}
    \item Sélectionnez le mode \textbf{Manuel}
    \item Chargez UN document PDF avec \bouton{Charger PDF}
    \item Dans le champ \champ{Contenu du QR code}, saisissez le texte souhaité :
    \begin{itemize}
        \item URL : \texttt{https://www.fmsp.edu/verify?id=123}
        \item Texte : \texttt{Document officiel FMSP 2024}
        \item vCard : Contact au format vCard
        \item Tout autre texte libre
    \end{itemize}
    \item Configurez la position avec la grille A4
    \item Choisissez la taille et la correction d'erreur
    \item Cliquez sur \bouton{Générer le QR code}
    \item Un aperçu s'affiche avec le QR code positionné
    \item Si satisfait, cliquez sur \bouton{Télécharger le PDF}
\end{enumerate}

\begin{figure}[h]
    \centering
    % \includegraphics[width=0.9\textwidth]{images/qr_placement_manuel.png}
    \fbox{\parbox{0.85\textwidth}{\centering [CAPTURE D'ÉCRAN : MODE MANUEL]\\ \vspace{0.3cm}
    Champ texte "Contenu du QR code"\\
    Grille de positionnement\\
    Aperçu du PDF avec QR code}}
    \caption{Interface du mode manuel}
\end{figure}

\subsection{Mode 2 : Traitement en lot avec Excel}

\subsubsection{Cas d'usage}

Utilisez ce mode pour :
\begin{itemize}[leftmargin=*]
    \item Traiter plusieurs documents en une seule opération
    \item Associer automatiquement des données Excel à chaque document
    \item Générer des QR codes uniques basés sur les données de chaque étudiant
\end{itemize}

\subsubsection{Workflow Excel}

\begin{enumerate}[leftmargin=*]
    \item Sélectionnez le mode \textbf{Excel}
    \item Chargez plusieurs PDFs avec \bouton{Charger Documents PDF}
    \item Chargez le fichier Excel avec \bouton{Charger Fichier Excel}
    \item Sélectionnez la \champ{Colonne de correspondance} :
    \begin{itemize}
        \item MATRICULE (recommandé)
        \item NOM
        \item Toute autre colonne unique
    \end{itemize}
    \item L'application tente de matcher automatiquement chaque PDF
    \item Corrigez manuellement les correspondances non trouvées
    \item Sélectionnez les colonnes à inclure dans les QR codes
    \item Configurez le positionnement et les paramètres du QR
    \item Choisissez le format d'export (ZIP / Individual / Single)
    \item Cliquez sur \bouton{Traiter tous les documents}
\end{enumerate}

\subsection{Paramètres avancés du QR code}

\subsubsection{Taille personnalisée}

\begin{enumerate}[leftmargin=*]
    \item Sélectionnez \champ{Taille : Personnalisée}
    \item Un champ s'affiche : \champ{Taille en pixels}
    \item Entrez une valeur entre 40 et 200 pixels
    \item Recommandations :
    \begin{itemize}
        \item 40-60 px : Très discret, pour documents chargés
        \item 80-100 px : Standard, bon compromis
        \item 120-150 px : Grand, très visible
        \item 150+ px : Très grand, pour affiches
    \end{itemize}
\end{enumerate}

\subsubsection{Correction d'erreur}

Les 4 niveaux expliqués en détail :

\begin{description}[leftmargin=2.5cm, style=nextline]
    \item[L (Low)]
    \begin{itemize}
        \item 7\% de données récupérables si endommagées
        \item QR code plus simple (moins de modules)
        \item Idéal pour : documents imprimés sur imprimante de qualité, environnement propre
        \item Éviter si : photocopies multiples, scan de mauvaise qualité
    \end{itemize}

    \item[M (Medium)]
    \begin{itemize}
        \item 15\% de données récupérables
        \item Bon équilibre entre densité et robustesse
        \item \textbf{Recommandé pour la plupart des cas}
        \item Usage type : relevés de notes, attestations standard
    \end{itemize}

    \item[Q (Quartile)]
    \begin{itemize}
        \item 25\% de données récupérables
        \item QR code plus dense
        \item Idéal pour : documents qui seront photocopiés, archivage long terme
        \item Usage type : diplômes, certificats officiels
    \end{itemize}

    \item[H (High)]
    \begin{itemize}
        \item 30\% de données récupérables
        \item QR code très dense
        \item Idéal pour : environnements difficiles, affichage public, documents manipulés fréquemment
        \item Usage type : affiches, documents d'identité
    \end{itemize}
\end{description}

\subsection{Prévisualisation avant traitement}

Avant de lancer un traitement en masse :

\begin{enumerate}[leftmargin=*]
    \item Configurez tous les paramètres
    \item Cliquez sur \bouton{Prévisualiser avec le premier document}
    \item Un PDF de test s'affiche avec le QR code positionné
    \item Vérifiez :
    \begin{itemize}
        \item Position correcte
        \item Taille appropriée
        \item Lisibilité du QR code (testez avec un scanner)
        \item Absence de chevauchement avec du texte important
    \end{itemize}
    \item Si nécessaire, ajustez les paramètres et prévisualisez à nouveau
    \item Une fois satisfait, lancez le traitement complet
\end{enumerate}

\subsection{Gestion des correspondances}

\subsubsection{Matching automatique}

L'algorithme de matching utilise plusieurs stratégies :

\begin{enumerate}[leftmargin=*]
    \item \textbf{Recherche exacte} : Le nom du fichier contient la valeur exacte de la colonne de correspondance
    \item \textbf{Recherche partielle} : Recherche insensible à la casse et aux accents
    \item \textbf{Extraction de mots-clés} : Si le fichier est nommé "Releve\_2023001\_DUPONT.pdf", extraction de "2023001" et "DUPONT"
\end{enumerate}

\subsubsection{Correction manuelle}

Pour les documents non matchés automatiquement :

\begin{enumerate}[leftmargin=*]
    \item Le document apparaît en \textcolor{red}{rouge} avec statut "Non matché"
    \item Cliquez sur \bouton{Sélectionner manuellement}
    \item Une liste déroulante s'affiche avec tous les étudiants de l'Excel
    \item Utilisez la barre de recherche pour trouver rapidement
    \item Sélectionnez l'étudiant correspondant
    \item Le statut passe à \textcolor{green}{✓ Matché}
\end{enumerate}

\subsection{Formats d'exportation}

Les mêmes 3 formats que les autres fonctionnalités :

\begin{description}[leftmargin=3cm]
    \item[ZIP]
    \begin{itemize}
        \item Tous les documents dans une archive
        \item Nom : \texttt{Documents\_QR\_[Date]\_[Heure].zip}
    \end{itemize}

    \item[Fichiers individuels]
    \begin{itemize}
        \item Chaque document téléchargé séparément
        \item Délai de 100 ms entre chaque téléchargement
        \item Nom conservé : \texttt{[Nom\_original]\_avec\_QR.pdf}
    \end{itemize}

    \item[PDF unique]
    \begin{itemize}
        \item Tous les documents fusionnés
        \item Nom : \texttt{Documents\_Fusionnes\_QR\_[Date].pdf}
        \item Ordre : alphabétique par nom de fichier
    \end{itemize}
\end{description}

\newpage
